% sec:spec-lgl — LGL Differentiation and Mass Lumping
\section{LGL Differentiation and Mass Lumping}
\label{sec:spec-lgl}

The DG solver doesn't use a spectral differentiation matrix in the
same way as the collocation solver of \cref{sec:spec-diffmat}.
Collocation methods apply the $D$ matrix directly to function values;
DG methods work in weak form, multiplying by test functions and
integrating.  The weak-form machinery requires three matrices
assembled from the LGL nodes and weights of
\cref{sec:spec-legendre}: a differentiation matrix, a mass matrix,
and a stiffness matrix.  Together with the Vandermonde matrix
(\cref{sec:spec-vandermonde}) for modal--nodal conversion, all four
live inside the \texttt{LglBasis} struct that the codebase constructs
once at startup and reuses at every time step.
\coderef{dg}{LglBasis}

\paragraph{The LGL differentiation matrix.}

The LGL differentiation matrix is constructed by the same recipe
as the CGL version (\cref{sec:spec-diffmat}): the entry $D_{ij} =
\ell_j'(x_i)$ records the derivative of the $j$-th Lagrange
interpolant evaluated at the $i$-th node.  What changes is the node
set --- LGL nodes instead of CGL --- and consequently the formula for
the off-diagonal entries.  On CGL nodes the closed-form entry uses
endpoint flags $c_i$ (\cref{eq:diffmat-cgl}); on LGL nodes the
analogous formula uses evaluations of the Legendre polynomial $P_N$:
%
\begin{equation}
  D_{ij} = \frac{P_N(x_i)}{P_N(x_j)}
    \cdot \frac{1}{x_i - x_j} \,,
  \qquad i \neq j \,.
  \label{eq:lgl-diffmat-offdiag}
\end{equation}
%
The values $P_N(x_i)$ are already available from the Newton iteration
that computed the nodes (\cref{sec:spec-legendre}), so the matrix
construction requires no additional polynomial evaluations.  The
diagonal entries follow from the same row-sum property as before
(\cref{eq:diffmat-rowsum}): setting $D_{ii} = -\sum_{j \neq i}
D_{ij}$ guarantees that the matrix differentiates constants exactly.
The two endpoint entries have closed forms that bypass the row sum:
%
\begin{equation}
  D_{00} = -\frac{N(N+1)}{4} \,, \qquad
  D_{NN} = \frac{N(N+1)}{4} \,,
  \label{eq:lgl-diffmat-endpoints}
\end{equation}
%
with all interior diagonals computed from the negative row sum.
\implnote{The codebase function \texttt{lgl\_diff\_matrix} in
\codemodule{dg} implements \cref{eq:lgl-diffmat-offdiag} with the
endpoint formulas \cref{eq:lgl-diffmat-endpoints}, then fills
interior diagonals by negative row summation --- identical in
structure to the CGL code in \codemodule{spectral}.}

\paragraph{The mass matrix.}

The DG weak form multiplies the governing PDE by a test function
$\phi_i$ and integrates over a single element.  In the Lagrange nodal
basis, the test functions are the Lagrange interpolants $\ell_i$
themselves, and the mass matrix that results from the left-hand side
has entries
%
\begin{equation}
  M_{ij} = \int_{-1}^{1} \ell_i(x)\,\ell_j(x)\,\d x \,.
  \label{eq:mass-matrix}
\end{equation}
%
This integral involves a product of two degree-$N$ polynomials ---
an integrand of degree $2N$.  The LGL quadrature rule with $N + 1$
nodes is exact for polynomials of degree at most $2N - 1$
(\cref{sec:spec-legendre}), so it misses the mark by one degree.
Approximating the integral by the LGL rule gives the
\emph{lumped mass matrix}
%
\begin{equation}
  \hat{M}_{ij}
  = \sum_{k=0}^{N} w_k\,\ell_i(x_k)\,\ell_j(x_k)
  = w_i\,\delta_{ij} \,,
  \label{eq:mass-lumped}
\end{equation}
%
where the last step uses $\ell_i(x_k) = \delta_{ik}$ (the cardinal
property of Lagrange interpolants at the collocation nodes).  The
lumped mass matrix is diagonal: its only entries are the LGL
quadrature weights $w_i$ from \cref{eq:lgl-weights}.  Inverting it
requires only $N + 1$ divisions, not a matrix solve.
\physnote{Mass lumping is exact for all but the highest mode.  The
quadrature error affects only the $(P_N, P_N)$ entry in the modal
basis, as established in \cref{sec:spec-legendre}.  For smooth
solutions, the highest mode carries exponentially small amplitude,
so the lumping error is negligible.}

\paragraph{Why mass lumping works.}

The one-degree shortfall deserves scrutiny: if the quadrature doesn't
integrate $\ell_i\,\ell_j$ exactly, why trust the result?  The
argument from \cref{sec:spec-legendre} showed that the Legendre modal
mass matrix is diagonal by orthogonality, and that the LGL quadrature
reproduces this diagonality for all entries except $(P_N, P_N)$.
In the nodal basis the same conclusion holds by a different route.
The Lagrange interpolants at LGL nodes satisfy an exact
discrete orthogonality --- the identity \cref{eq:mass-lumped}
holds because of the cardinal property, regardless of the
quadrature's polynomial exactness.  The lumped mass matrix is the
identity (up to the weight factors) by construction --- it is the
\emph{exact} integral of $\ell_i\,\ell_j$ that differs from it, and
the difference is confined to a single mode.  Mass lumping doesn't
introduce error by under-resolving a difficult integral; it
\emph{defines} the inner product that the DG scheme uses, and that
inner product is spectrally close to the exact one.
\warnnote{Mass lumping is not free: it modifies the numerical scheme
by replacing the exact $L^2$ inner product with the discrete LGL
inner product.  For $p$-refinement studies comparing different $N$
at fixed element count, the mass-lumping error decreases
exponentially with $N$ and is invisible in practice.  For very
coarse resolutions ($N = 1$ or $2$), the error can be noticeable.}

\paragraph{The stiffness matrix.}

The right-hand side of the DG weak form produces the stiffness
matrix, which encodes differentiation weighted by the test function.
Applying LGL quadrature to the integral
$\int \ell_i(x)\,\ell_j'(x)\,\d x$ --- an integrand of degree
$2N - 1$, so the quadrature is exact --- gives
%
\begin{equation}
  S_{ij} = \sum_{k=0}^{N} w_k\,\ell_i(x_k)\,\ell_j'(x_k)
  = w_i\,D_{ij} \,,
  \label{eq:stiffness-matrix}
\end{equation}
%
using $\ell_i(x_k) = \delta_{ik}$ and the differentiation matrix
definition $D_{ij} = \ell_j'(x_i)$ from \cref{eq:diffmat-def}.
The stiffness matrix is the mass-weighted differentiation matrix:
$S = \operatorname{diag}(w_0, \ldots, w_N) \cdot D$.  The
combination $\hat{M}^{-1} S = D$ --- the lumped mass inverse times
the stiffness matrix is just the differentiation matrix itself.
This identity is the algebraic payoff of mass lumping: the weak-form
DG operator reduces to a simple matrix-vector multiply by $D$,
with no matrix inversion at each time step.
\xrefnote{The DG solver applies $D$ in each coordinate direction
independently, producing the volume derivative terms.  Face
contributions from the numerical flux are added separately
(\cref{ch:dg-methods}).}

\paragraph{Assembly in the codebase.}

The \texttt{lgl\_basis} constructor in \codemodule{dg} assembles all
four components in sequence: nodes and weights from
\texttt{lgl\_nodes\_weights} (\cref{sec:spec-legendre}), the
differentiation matrix from \texttt{lgl\_diff\_matrix}
(\cref{eq:lgl-diffmat-offdiag}), and the Vandermonde matrix from
\texttt{legendre\_all} (\cref{sec:spec-vandermonde}).  The resulting
\texttt{LglBasis} struct stores:
%
\begin{lstlisting}[language=Rust]
pub struct LglBasis {
    pub order: usize,          // polynomial order p
    pub n: usize,              // number of nodes: p+1
    pub nodes: Vec<f64>,       // LGL nodes in [-1,1]
    pub weights: Vec<f64>,     // LGL quadrature weights
    pub d_matrix: Vec<f64>,    // D, n*n row-major
    pub v_matrix: Vec<f64>,    // Vandermonde V (modal->nodal)
    pub v_inv: Vec<f64>,       // V^{-1}
}
\end{lstlisting}
%
The struct is constructed once per simulation and shared across all
elements of the same polynomial order.  At each Runge--Kutta stage,
the volume derivative $D\,\mathbf{u}$ is computed by a dense
matrix-vector multiply; the mass matrix never appears explicitly
because the lumped inverse is absorbed into the operator.  For the
polynomial orders used in practice ($N = 3$--$7$), the $D$-matrix
multiply dominates the cost of each element update, and its
$\order{N^2}$ operation count per node is modest --- a single
element with $N = 5$ requires $36$ multiplies per derivative
direction.
\implnote{In three dimensions each element stores
$(N+1)^3$ unknowns per field variable.  The volume derivative
applies $D$ along each axis independently via tensor-product
structure, reducing the per-element cost from
$\order{N^6}$ (dense matrix) to $\order{N^4}$ (three
one-dimensional passes) --- the sum-factorisation technique of
\cite{Kopriva2009}.}
